\chapter{Limitations, Risks, and Evolution}
\label{cap:limitacoes}

\section{Technical limitations}

\begin{table}[H]
\centering
\caption{Limitations by area and consequences.}
\label{tab:limitacoes}
\begin{tabularx}{\linewidth}{@{}p{3.3cm}X@{}}
\toprule
\textbf{Area} & \textbf{Limitation and consequence} \\
\midrule
Interrupts & IDT with 256 gates, remapped PIC, and IRQ0 timer. The keyboard's
IRQ1 remains masked; there is no APIC or SMP, and IRQ0 drives the preemptive
scheduler. \\
Memory & Frame allocator, 4 KiB pages on demand, and a 1 MiB first-fit kernel
heap with \code{Vec}/\code{String}. No slab or buddy allocator, no swap, no
reclaim, and no NUMA. \\
Processes & Task table, PID/PPID, context switch, a timer-preemptive
round-robin scheduler with a five-tick quantum, blocking \code{sleep}
(\code{schedule\_timeout}) and event wait queues
(\code{wait\_event}/\code{wake\_up}). There is no
\code{fork}/\code{exec}/\code{wait}, no CFS/EEVDF, no priorities and no
per-process address space; the ring 3 programs are still not scheduled
processes. Eleven syscalls
(\code{exit/write/read/get\_tty/open/close/read\_fd/write\_fd/lseek/getpid/
getppid}). \\
Storage & VFS with sixteen descriptors over the in-memory ramfs and a synthetic
\code{/proc}. No disk driver, block layer, mounts, permissions, page cache,
FAT, or ext2; the ISO is only the boot medium. \\
Networking & No TCP/IP stack, netdevice, or network drivers. \\
Video & Legacy VGA text only. UEFI without CSM, GPUs, and graphics
modes are not supported. \\
Input & PS/2 keyboard by polling, with no native USB-HID, Caps Lock, AltGr,
arrow keys, or ABNT2 layout. \\
Serial & COM1 only, polling, no IRQ, DMA, or flow control. \\
Concurrency & One CPU; the scheduler is timer-preemptive with a five-tick
quantum, a blocking \code{sleep} and event wait queues, but preemption is
suspended while a ring 3 program is active. The spinlocks use \code{irqsave},
so the task table is safe against IRQ0. \\
Security & U/S separates user and kernel, but there is no ASLR, NX,
SMEP/SMAP, \emph{stack canary}, or per-process address spaces. User pointers
are range-checked, but there is no \code{copy\_from\_user}-style helper. \\
Robustness & The IDT prevents \emph{triple fault} in most cases; there is no
failure recovery, watchdog, or persistent logs. \\
Portability & Code specific to x86\_64 BIOS; it does not compile for
ARM, RISC-V, or pure UEFI. \\
\bottomrule
\end{tabularx}
\end{table}

\section{Risks}

\begin{table}[H]
\centering
\caption{Risks and current mitigations.}
\label{tab:riscos}
\begin{tabularx}{\linewidth}{@{}p{3.7cm}p{3.0cm}X@{}}
\toprule
\textbf{Risk} & \textbf{Likelihood} & \textbf{Mitigation} \\
\midrule
\emph{Triple fault} due to an exception & Medium for new code & IDT and
\emph{stubs} already installed; test new handlers in QEMU before
enabling critical paths. \\
\emph{Deadlock} in the console & Low today & Do not acquire the lock twice;
the timer does not print, but future IRQs must use \code{irqsave}. \\
\emph{Hang} while waiting for hardware & Medium & Timeouts in the 8042 waits;
the UART can still hang if present and faulty. \\
Keyboard does not work on real hardware & Medium & Test with \emph{legacy
USB}/PS/2; use serial as an alternative. \\
Black screen on UEFI & High without CSM & Document the BIOS/CSM requirement;
implement GOP in the future. \\
Byte loss on serial & Medium & Fast polling; FIFO/IRQ in the future. \\
Corruption due to an index bug & Low & Review the invariants of
\code{row}/\code{column}; add tests. \\
\bottomrule
\end{tabularx}
\end{table}

\section{Suggested roadmap}

The order below prioritizes gains in robustness and understanding before
adding complexity.

\begin{table}[H]
\centering
\caption{Prioritized roadmap.}
\label{tab:roadmap}
\begin{tabularx}{\linewidth}{@{}clX@{}}
\toprule
\textbf{Prior.} & \textbf{Item} & \textbf{Goal} \\
\midrule
1 & Priorities and signals & Weighted scheduling classes and signal delivery
(plus per-task kernel stacks for scheduled ring 3 tasks). \\
2 & Real processes & \code{fork}/\code{exec}/\code{wait}, per-process page
tables, and ring 3 tasks scheduled by the kernel. \\
3 & VFS and filesystem & Inodes, dentries, mounts and permissions; a disk
driver and an initramfs. \\
4 & IRQ and keyboard robustness & IRQ1 with a scancode buffer and handling of
spurious IRQs. \\
5 & TTY line discipline & Canonical mode, signals, job control,
\code{Ctrl+C}/\code{Ctrl+D} and one \code{getty} per console. \\
6 & GOP/VBE framebuffer & Support for UEFI and graphics modes; bitmap font. \\
7 & Memory hardening & NX, SMEP/SMAP, validation of user
pointers, and ASLR. \\
8 & Complete keyboard driver & USB-HID, Caps Lock, AltGr, ABNT2 layout,
extended keys. \\
9 & Networking & NE2000/virtio-net, ARP/IP/ICMP/UDP, tests with QEMU. \\
10 & Automated tests & CI with headless QEMU, VGA dumps, and asserts. \\
\bottomrule
\end{tabularx}
\end{table}

\section{Criteria for considering the project ``a kernel''}

From the educational point of view, the project is already a kernel: it runs in ring 0,
controls hardware, and does not depend on an operating system. From the point of view
of a usable operating system, the following are still missing, at a minimum:

\begin{itemize}
  \item process isolation and a per-process user space;
  \item scheduling priorities, signals and per-task kernel stacks for ring 3
        tasks (blocking sleep and event wait queues already exist);
  \item a persistent filesystem and a program \emph{loader};
  \item input/output drivers beyond VGA/PS/2/serial;
  \item error handling and recovery.
\end{itemize}

The current project is best classified as a \textbf{minimal monolithic
demonstration kernel}, focused on boot, memory, interrupts, files, tasks,
and console.
